\section{Aprendizaje continuo y colaboración multi-agent}

\subsection{Manifestación multi-agent e imitation}

\begin{figure}[H]
\centering
\includegraphics[width=0.8\textwidth]{slide-02-12.jpg}
\caption{Slide 02-12 muestra el aprendizaje por imitación y los communication loops en un multi-agent scenario.}
\end{figure}

En un escenario multi-agent, cada agent puede imitar a distintos specialist, y comparten policy logits entre sí mediante una gating network. La CSI gating permite elegir, en runtime, el expert más adecuado.

\begin{importantbox}{El poder de multi-agent imitation}
Mediante multi-agent imitation, el sistema puede dividir el trabajo y colaborar en tareas complejas que un modelo de un solo agent no puede cubrir. Cada agent se encarga de un segmento de trajectory y, al final, un aggregator sintetiza una política consistente.
\end{importantbox}

\subsection{Aprendizaje continuo y automejora}

\begin{figure}[H]
\centering
\includegraphics[width=0.8\textwidth]{slide-02-13.jpg}
\caption{Slide 02-13 dibuja el continuous improvement loop: data → model → evaluation → deployment.}
\end{figure}

Dentro del continuous improvement loop, cada deployment debe tomar el human override como new training data, e incorporar un multi-agent critic como auxiliary loss dentro del nightly retrain pipeline.

\begin{knowledgebox}{Self-supervised critic}
Usar un self-supervised critic permite estimar el reward gap cuando falta human comparison, y combinado con human override forma una stratified preference signal.
\end{knowledgebox}

\subsection{Resumen de la sección}

Multi-agent imitation y el continuous improvement loop convierten el aprendizaje por imitación en un sistema que se autocalibra constantemente; basta con agregar traction guardrails en evaluation → retrain → deployment para lograr una mejora continua.

